```latex
\chapter*{Conclusion}
\addcontentsline{toc}{part}{Conclusion}

\section*{Summary and conclusion}

The main objective of this thesis was to investigate the possibilities of electronic structure engineering in group IV crystals, with particular emphasis on materials compatible with the existing silicon-based technology and potentially suitable for optoelectronic applications. Four approaches to modifying the electronic structure were considered: strain, alloying, structural modifications, and quantum confinement. The calculations were performed within the framework of density functional theory, using different computational approaches and exchange--correlation functionals depending on the properties under investigation.

The results demonstrate that even relatively simple modifications of group IV crystal structures can lead to substantial changes in their electronic properties. In particular, the fundamental band gap, its direct or indirect character, band offsets, and effective masses can be strongly affected. At the same time, the investigations provide structural and elastic parameters that are relevant for understanding and predicting the response of these materials to external or internally induced modifications.

\subsection*{Chapter~\ref{chap:strains} -- Strain engineering}

The aim of this part of the work was to provide a systematic description of the influence of technologically relevant strains on the electronic structures of group IV crystals in the diamond structure. Isotropic, uniaxial, and biaxial strains were investigated, with the latter two applied along the crystallographically relevant $[100]$, $[110]$, and $[111]$ directions.

The calculations have shown that strain produces significant changes in the band structures, particularly in the fundamental band gaps and the relative positions of the band extrema. Depending on the material, strain type, and crystallographic direction, transitions between direct and indirect band gaps can occur. The symmetry reduction associated with non-isotropic strains also leads to the splitting of formerly degenerate states at selected high-symmetry points of the Brillouin zone. These effects demonstrate that the electronic response cannot be described by considering only the magnitude of strain; its geometrical character and direction are equally important.

The deformation potentials obtained in this work provide a quantitative description of the changes in band-edge energies and can be used in further modelling of strained semiconductor structures. The effective masses of electrons and holes were also evaluated at relevant band extrema. Particularly important changes were observed in the valence band, where strain-induced changes in the ordering of the heavy-hole and light-hole bands may result in sudden variations of the corresponding effective masses. In some cases, these changes reach approximately one order of magnitude. The effective mass of electrons in Ge near the $\Gamma$ point also exhibits a particularly strong strain dependence.

The estimated carrier mobilities reflect these changes in effective masses. Under the assumption that the scattering time remains unchanged, the calculated mobility can vary considerably with strain, in extreme cases by approximately one order of magnitude. This indicates that strain engineering may be used not only to modify the optical properties of group IV materials but also to influence their charge transport characteristics.

An additional outcome of this investigation is a consistent set of structural and elastic parameters, including lattice constants, elastic coefficients, Poisson's ratios, and strain-induced stresses. The comparison with available experimental and theoretical data provides support for the reliability of the adopted computational procedures. These results constitute a useful reference for further studies of strained group IV materials.

\ornament

\subsection*{Chapter~\ref{chap:alloying} -- Alloying}

The alloying part of the work focused on the application of the computational alchemy approximation to the electronic structure of GeSn alloys. The main objective was to reproduce the electronic band structures obtained from supercell calculations while retaining the smooth band dispersion and computational efficiency of the alchemical approach.

The lattice and $c_{\mathrm{TB09}}$ parameters were adjusted to obtain agreement between the calculated and reference supercell results. The comparison demonstrated very good consistency between the two approaches. In the investigated composition range, the differences in the direct and indirect band gaps did not exceed approximately $0.03$ eV. A similarly good agreement was obtained for the band structures and electron effective masses at the $\Gamma$ and $L$ valleys.

The results show that the computational alchemy approximation can provide a useful representation of the electronic structure of GeSn alloys, avoiding some of the computational complexity associated with explicit supercell calculations. The smooth band structures obtained in this way are particularly suitable for further modelling of semiconductor systems. Indeed, the results of this investigation were subsequently used in the development of a 30-band $\mathbf{k}\cdot\mathbf{p}$ model for GeSn, demonstrating the practical usefulness of the approach.

\ornament

\subsection*{Chapter~\ref{chap:wurtzite} -- Structural factor}

The influence of crystal structure on the physical properties of group IV materials was investigated by considering the wurtzite phase of C, Si, and Ge. The structural and elastic properties were determined first, including lattice constants, the internal structural parameter $u$, bulk moduli, elastic constants, pressure coefficients, and Poisson's ratios. The formation energies of the wurtzite and zinc blende structures were also compared under different strain conditions.

The results show that the general chemical trends observed in the diamond structure are largely preserved in the wurtzite phase. In particular, the elastic stiffness decreases when moving from carbon towards the heavier group IV elements. The comparison of formation energies indicates that the zinc blende structure is generally energetically preferred. An exception is Sn under isotropic tensile strain, for which the wurtzite structure becomes energetically favourable. Nevertheless, the energy differences are relatively small, suggesting that the energetic criterion alone may not be sufficient to determine the experimentally accessible structure.

The electronic structure calculations revealed particularly interesting properties of wurtzite Ge. Unlike Si and C, Ge exhibits a direct fundamental band gap in its equilibrium configuration. Furthermore, the direct gap is preserved over a broad range of applied strains, except for sufficiently strong biaxial compression. Its magnitude can be modified substantially, ranging from an almost closed gap to values exceeding $0.6$ eV under isotropic compression. Together with the calculated electron and hole effective masses below $0.2\,m_e$ along the $\Delta$ direction, these results indicate the potential of wurtzite Ge for infrared optoelectronic applications.

For Si, a transition to a direct fundamental band gap was also found under biaxial tensile strain, although it requires strain exceeding approximately $4\%$. The calculated spin--orbit and crystal-field splittings provide additional information about the electronic structure of these systems. While the spin--orbit splitting changes only weakly with strain, the crystal-field splitting exhibits a considerably stronger response.

Overall, the results demonstrate that changing the crystal structure, combined with strain engineering, offers another route towards modifying the electronic properties of group IV materials. In particular, wurtzite Ge appears to be a promising system for further investigation.

\ornament

\subsection*{Chapter~\ref{chap:digital} -- Spatial confinement}

The final part of the work provided a preliminary investigation of Ge--Si digital alloys, in which the constituent materials form an ordered layered structure. Systems grown along the $[100]$ direction were considered on both Si and Ge substrates. The structural calculations were verified using the Poisson's ratio approach, which produced atomic positions and lattice parameters consistent with those obtained from structural relaxation.

The electronic band structures were analysed along a Brillouin zone path covering both in-plane and out-of-plane directions. The atom-projected band structures provided additional insight into the spatial character of the electronic states and allowed the contributions of individual atomic layers to be distinguished.

The calculations revealed differences between systems grown on Si and Ge substrates. In both cases, band folding and direct band-gap features of the order of $1$ eV were observed. However, the fundamental gap and the character of the electronic states depend on the substrate and on the atomic contributions considered. In particular, for the system on a Si substrate, the atom-projected band structures indicate a direct gap for selected atomic contributions, despite the indirect character of the fundamental gap in the total band structure. For the Ge substrate, a direct gap is observed around the $\Gamma$ point, with additional characteristic features in the dispersion perpendicular to the layers.

These observations suggest that Ge--Si digital alloys may exhibit spatially dependent electronic properties resembling those of quantum-well systems. The results are preliminary, but they demonstrate that the layered arrangement of group IV materials provides an additional degree of freedom for engineering their electronic structures. They also establish a basis for further investigations of layer thickness, growth direction, spatially resolved band gaps, and optical transitions.

\ornament

\section*{Achieved goals}

The investigations presented in this thesis have addressed the four principal approaches to electronic structure engineering introduced in the first part of the work. The main goals and their outcomes can be summarised as follows:

\begin{enumerate}
	\item A systematic study of strain effects in group IV crystals in the diamond structure was performed. Different types and crystallographic directions of strain were considered, and their influence on structural, elastic, and electronic properties was evaluated. The resulting data provide a consistent description of strain-dependent band gaps, deformation potentials, effective masses, and estimated carrier mobilities.
	
	\item The computational alchemy approximation was applied to GeSn alloys and validated against reference supercell calculations. The good agreement obtained for band structures, band gaps, and effective masses confirms that this approach can be used as an efficient tool for describing alloy electronic structures and for generating input data for further semiconductor modelling.
	
	\item The structural and electronic properties of group IV materials in the wurtzite phase were investigated. In particular, the direct band gap of equilibrium wurtzite Ge was identified, and its response to different types of strain was evaluated. The results establish a theoretical basis for considering this material in the context of infrared optoelectronic applications.
	
	\item The electronic structure of Ge--Si digital alloys grown along the $[100]$ direction was investigated for different substrate materials. The atom-projected band structures revealed differences in the spatial character of electronic states and indicated the possibility of spatially dependent band-gap properties in layered systems.
	
	\item The computational procedures and results were compared with available experimental and theoretical reference data wherever possible. This included structural parameters, band structures, band gaps, and effective masses. Such comparisons, together with convergence tests and the use of complementary computational approaches, provided an assessment of the reliability of the calculated properties.
\end{enumerate}

Taken together, these results show that the electronic properties of group IV materials can be modified through several complementary approaches. Strain and structural modifications can change the character and magnitude of the fundamental band gap, alloying provides a means of continuously tuning the electronic structure, while layered systems introduce additional spatial degrees of freedom. The work therefore contributes both numerical data and methodological approaches that may be useful in the further design and modelling of group IV semiconductor systems.

\ornament

\section*{Challenges}

Despite the progress achieved, several challenges and limitations remain. They concern both the computational methods and the physical interpretation of the results.

One of the main difficulties in first-principles calculations is the balance between computational cost and the accuracy required for a given physical property. Different exchange--correlation functionals and pseudopotentials were therefore used depending on whether structural or electronic properties were investigated. Although the adopted approaches produced satisfactory agreement with available reference data, the accuracy of individual quantities is not uniform. In particular, effective masses calculated using the modified Becke--Johnson potential may contain systematic errors, and their quantitative values should be interpreted with appropriate caution. The results provide a reliable qualitative description of their strain dependence, but further validation would be desirable.

The determination of deformation potentials is another demanding task. While the adopted procedure is based on a well-defined theoretical framework, reference data for some of the investigated systems are limited or unavailable. This makes an independent assessment of the absolute values difficult. Further calculations using alternative computational approaches, as well as comparison with experimental measurements, would help to establish their quantitative reliability.

The mobility estimates also involve an important simplifying assumption. The strain dependence was evaluated through changes in the effective masses while assuming that the carrier scattering time remains unaffected by strain. In real materials, however, strain may also modify scattering mechanisms and therefore the scattering time itself. A more complete description would require explicit consideration of the relevant scattering processes, which is beyond the scope of the present work.

In the alloying study, the computational alchemy approximation was validated for the investigated GeSn composition range and against the available supercell reference data. Its applicability to other alloy compositions, different material systems, or more complex structural configurations requires further verification. In particular, the range of compositions and the treatment of possible structural and local-environment effects could be extended.

The wurtzite investigation is limited by the scarcity of experimental reference data for the considered group IV materials. Although the calculated properties of Si show satisfactory agreement with available theoretical results, further experimental or independent theoretical validation would be valuable. Moreover, the energetic preference of a crystal structure does not by itself determine whether a given phase can be stabilised experimentally. The role of growth conditions, substrates, and kinetic effects remains to be investigated.

Finally, the study of Ge--Si digital alloys should be regarded as preliminary. The present calculations considered a limited set of layered structures and one principal growth direction. The spatial distribution of the band edges has not yet been determined in sufficient detail to establish the full quantum-well character of the systems. Further work should include different layer thicknesses, growth along the $[111]$ direction, spatially resolved electronic states, and a detailed analysis of optical transitions and oscillator strengths. The extension of this approach to wurtzite digital alloys is another promising direction.

These challenges define natural directions for further research. In particular, combining the structural and electronic information obtained in this thesis with more detailed optical and transport calculations would allow a more direct assessment of the suitability of the investigated systems for optoelectronic applications.

\ornament

In conclusion, this thesis demonstrates that group IV materials offer several complementary possibilities for electronic structure engineering. Although the investigated systems differ in their structural complexity and technological maturity, the results consistently show that their electronic properties can be tailored through strain, alloying, crystal structure modifications, and spatial confinement. The findings provide a basis for further theoretical and experimental studies aimed at extending the functionality of silicon-compatible materials towards optoelectronic applications.
```